\documentclass[12pt,a4paper]{article}

% Packages
\usepackage[utf8]{inputenc}
\usepackage[margin=1in]{geometry}
\usepackage{graphicx}
\usepackage{amsmath}
\usepackage{amssymb}
\usepackage{booktabs}
\usepackage{hyperref}
\usepackage{caption}
\usepackage{subcaption}
\usepackage{float}
\usepackage{listings}
\usepackage{xcolor}
\usepackage{multirow}
\usepackage{enumitem}
\usepackage{multicol}

% Code listing style
\lstset{
    basicstyle=\ttfamily\small,
    breaklines=true,
    frame=single,
    backgroundcolor=\color{gray!10}
}

% Title
\title{\textbf{ITI123 Project Milestone Report} \\
\Large{AI-Based Badminton Coaching System: \\
Stroke Technique Assessment Using Deep Learning}}

\author{Student Name \\ Student ID \\ ITI123 - Artificial Intelligence}
\date{January 29, 2026}

\begin{document}

\maketitle

\begin{abstract}
This milestone report presents progress on developing an AI-based badminton coaching system that provides objective, quantitative feedback on stroke technique execution. The system focuses on two fundamental overhead strokes (Clear and Smash) and aims to assist players in improving technique through pose-based analysis. We implemented a comprehensive pipeline including MediaPipe pose extraction on the ShuttleSet dataset (4,983 clips, 40 matches), engineered 60 biomechanical features capturing key technique indicators (arm extension, wrist angles, body posture), and developed multiple model architectures. Through rigorous Cohen's d effect size analysis, we discovered that Clear and Smash strokes exhibit nearly identical body mechanics at shuttlecock contact (Cohen's d < 0.2), which has important implications for coaching systems. Rather than viewing this as a failure, we reframe these findings as valuable coaching insights: the primary differences between these strokes lie in racket control and shuttlecock trajectory, not body pose. This report documents our methodology, analytical findings, and proposes an alternative coaching-focused approach that leverages pose features to assess technique quality within each stroke type, rather than attempting to distinguish between them.
\end{abstract}

\tableofcontents
\newpage

\section{Introduction}

\subsection{Problem Statement and Business Context}

\textbf{Business Problem}: Badminton players, especially recreational and intermediate-level athletes, lack access to consistent, objective feedback on their stroke technique. Traditional coaching relies on expert observation, which is subjective, expensive, and difficult to scale. Players need automated systems that can analyze their technique and provide actionable feedback for improvement.

\textbf{Proposed Solution}: An AI-based coaching system that analyzes video clips of badminton strokes and provides quantitative assessment of technique quality. The system focuses on two fundamental overhead strokes:

\begin{itemize}
    \item \textbf{Clear}: Defensive stroke sending the shuttlecock high to the opponent's backcourt
    \item \textbf{Smash}: Aggressive attacking stroke with steep downward trajectory
\end{itemize}

\textbf{User Need}: Players require feedback on movement characteristics such as arm extension, body posture, timing, and swing mechanics to identify technical flaws and improve performance.

\textbf{Initial Research Objective}: As a first step, we investigated whether these two stroke types could be automatically distinguished using pose-only features, which would enable the system to provide stroke-specific coaching feedback.

\subsection{Dataset}

\textbf{Source}: ShuttleSet - A Human Action Recognition Dataset for Badminton Singles

\textbf{Statistics}:
\begin{itemize}
    \item Total clips: 4,983 (extracted at shuttlecock contact)
    \item Clear strokes: 2,660 (53.4\%)
    \item Smash strokes: 2,321 (46.6\%)
    \item Unique matches: 40
    \item Pose extraction success rate: 99.4\% (4,954/4,983)
\end{itemize}

\subsection{Methodology Overview}

Our approach follows a systematic pipeline aligned with coaching system development:
\begin{enumerate}
    \item \textbf{Pose extraction} using MediaPipe to capture body kinematics
    \item \textbf{Feature engineering} focused on technique indicators: arm extension, wrist angles, body posture, swing velocity
    \item \textbf{Data splitting} with match-level stratification to ensure generalization
    \item \textbf{Exploratory analysis} to understand biomechanical differences between strokes
    \item \textbf{Statistical validation} using Cohen's d effect size to quantify technique differences
    \item \textbf{Model development} to explore pattern recognition capabilities
    \item \textbf{Insight extraction} for coaching applications
\end{enumerate}

\textbf{Key Pivot}: Initial experiments revealed that stroke distinction is not feasible with pose-only data. We pivot toward using these insights to inform an alternative coaching approach focused on technique quality assessment within each stroke type.

\section{Technical Implementation}

\subsection{Pose Extraction}

\textbf{Framework}: Google MediaPipe Pose \\
\textbf{Output}: 33 body landmarks per frame (x, y, z, visibility) \\
\textbf{Processing}:

\begin{lstlisting}[language=Python, caption=MediaPipe Pose Extraction]
import mediapipe as mp

mp_pose = mp.solutions.pose
pose = mp_pose.Pose(
    static_image_mode=False,
    model_complexity=2,  # Highest accuracy
    smooth_landmarks=True,
    min_detection_confidence=0.5
)

# Process video frame-by-frame
for frame in video:
    results = pose.process(frame)
    landmarks = results.pose_landmarks  # 33 keypoints
\end{lstlisting}

\textbf{Key Landmarks}: Shoulders, elbows, wrists, hips, knees, ankles, nose

\subsection{Feature Engineering}

We implemented two iterations of feature engineering:

\subsubsection{Version 1 (Baseline Features)}

\textbf{Features per frame}: 42

\begin{table}[H]
\centering
\caption{Baseline Feature Categories}
\begin{tabular}{@{}ll@{}}
\toprule
\textbf{Category} & \textbf{Features} \\ \midrule
Arm Extension & Distance from shoulder to wrist (L/R) \\
Depth (Z-coordinate) & Wrist/elbow depth relative to shoulder/hip \\
Height (Y-coordinate) & Wrist/elbow height relative to shoulder/head \\
Joint Angles & Elbow angle, shoulder angle \\
Body Posture & Torso lean, shoulder rotation \\
Velocities & First derivative of all spatial features \\ \bottomrule
\end{tabular}
\end{table}

\subsubsection{Version 2 (Enhanced with Wrist Orientation)}

\textbf{Features per frame}: 60 \\
\textbf{New additions}: 9 wrist/forearm orientation features

\begin{table}[H]
\centering
\caption{Wrist Orientation Features (New)}
\begin{tabular}{@{}lp{8cm}@{}}
\toprule
\textbf{Feature} & \textbf{Description} \\ \midrule
\texttt{r\_forearm\_vertical\_angle} & Angle between forearm and vertical axis (key discriminator hypothesis) \\
\texttt{r\_forearm\_horizontal\_angle} & Lateral swing direction in X-Z plane \\
\texttt{r\_wrist\_elbow\_vertical} & Wrist height relative to elbow (flexion indicator) \\
\texttt{r\_wrist\_horizontal\_reach} & Forward extension distance \\
\texttt{r\_arm\_plane\_pronation} & Pronation/supination from plane normal \\
\texttt{r\_arm\_plane\_vertical} & Vertical component of arm plane \\
\texttt{r\_forearm\_pitch} & Elevation angle in sagittal plane \\
\texttt{l\_forearm\_vertical\_angle} & Left arm vertical angle \\
\texttt{l\_wrist\_elbow\_vertical} & Left wrist-elbow separation \\ \bottomrule
\end{tabular}
\end{table}

\textbf{Rationale}: Clear and Smash should differ in wrist angle at contact:
\begin{itemize}
    \item \textbf{Clear}: Wrist extended, racket angled upward (expected $\sim$120-150°)
    \item \textbf{Smash}: Wrist flexed, racket angled downward (expected $\sim$30-60°)
\end{itemize}

\textbf{Implementation}:
\begin{lstlisting}[language=Python, caption=Forearm Vertical Angle Calculation]
# Forearm vector (elbow to wrist)
forearm_vector = wrist - elbow
forearm_unit = forearm_vector / np.linalg.norm(forearm_vector)

# Vertical axis (downward in image coordinates)
vertical = np.array([0, 1, 0])

# Compute angle
cos_angle = np.dot(forearm_unit, vertical)
vertical_angle = np.arccos(np.clip(cos_angle, -1, 1))
vertical_angle_degrees = np.degrees(vertical_angle)
\end{lstlisting}

\subsection{Data Splitting}

\subsubsection{Group-Stratified Split}

To prevent data leakage and maintain class balance:

\begin{lstlisting}[language=Python, caption=Group-Stratified Splitting Strategy]
# Step 1: Calculate dominant class for each match
for match_id in unique_matches:
    clear_count = count_strokes(match_id, 'Clear')
    smash_count = count_strokes(match_id, 'Smash')
    dominant_class = 'Clear' if clear_count > smash_count else 'Smash'

# Step 2: Stratify matches by dominant class
clear_matches = matches_with_dominant('Clear')
smash_matches = matches_with_dominant('Smash')

# Step 3: Split each stratum
train_matches = 70% of (clear_matches + smash_matches)
val_matches = 15% of (clear_matches + smash_matches)
test_matches = 15% of (clear_matches + smash_matches)
\end{lstlisting}

\textbf{Benefits}:
\begin{itemize}
    \item Prevents data leakage (clips from same match stay together)
    \item Maintains class balance across splits
    \item Realistic evaluation (generalizes to new matches)
\end{itemize}

\textbf{Results}:
\begin{table}[H]
\centering
\caption{Data Split Statistics}
\begin{tabular}{@{}lrrrr@{}}
\toprule
\textbf{Split} & \textbf{Clips} & \textbf{Matches} & \textbf{\% Clear} & \textbf{\% Smash} \\ \midrule
Train & 3,347 & 27 & 53.22\% & 46.78\% \\
Validation & 687 & 5 & 52.03\% & 47.97\% \\
Test & 920 & 8 & 55.10\% & 44.90\% \\
\midrule
\textbf{Overall} & \textbf{4,954} & \textbf{40} & \textbf{53.40\%} & \textbf{46.60\%} \\ \bottomrule
\end{tabular}
\end{table}

All splits within $\pm$2\% of overall distribution.

\subsubsection{Padding Masking}

\textbf{Problem}: Padded frames pollute normalization statistics

\textbf{Solution}: Normalize only real frames, re-pad with zeros after normalization

\begin{lstlisting}[language=Python, caption=Proper Padding Masking]
# Collect only real frames for normalization
real_frames = []
for seq, length in zip(sequences, seq_lengths):
    real_frames.append(seq[:length])  # Exclude padding

# Compute normalization from real frames only
all_real = np.vstack(real_frames)
mean = np.mean(all_real, axis=0)
std = np.std(all_real, axis=0) + 1e-8

# Normalize and re-pad
for seq, length in zip(sequences, seq_lengths):
    normalized = (seq[:length] - mean) / std
    if length < 50:
        padding = np.zeros((50 - length, n_features))
        normalized = np.vstack([normalized, padding])

# Save sequence lengths for Masking layer
data = {'X': normalized, 'seq_lengths': seq_lengths}
\end{lstlisting}

\textbf{Model Integration}:
\begin{lstlisting}[language=Python, caption=LSTM with Masking Layer]
from tensorflow.keras.layers import Masking, LSTM

model = Sequential([
    Masking(mask_value=0.0, input_shape=(50, 60)),  # Ignore padding
    LSTM(128, return_sequences=True),
    Dropout(0.3),
    LSTM(64),
    Dense(1, activation='sigmoid')
])
\end{lstlisting}

\subsection{Model Architectures}

\subsubsection{Baseline Models}

\textbf{Input}: Statistical summary features (427 features)

\begin{enumerate}
    \item \textbf{Random Forest}
    \begin{itemize}
        \item n\_estimators = 100
        \item max\_depth = 20
        \item class\_weight = 'balanced'
    \end{itemize}

    \item \textbf{Support Vector Machine (SVM)}
    \begin{itemize}
        \item kernel = 'rbf'
        \item C = 1.0
        \item class\_weight = 'balanced'
        \item StandardScaler preprocessing
    \end{itemize}
\end{enumerate}

\subsubsection{Deep Learning Models}

\textbf{Input}: Sequence features (50 timesteps, 60 features)

\begin{enumerate}
    \item \textbf{LSTM}
    \begin{lstlisting}[language=Python]
Sequential([
    Masking(mask_value=0.0),
    LSTM(128, return_sequences=True),
    Dropout(0.3),
    LSTM(64),
    Dropout(0.3),
    Dense(1, activation='sigmoid')
])
    \end{lstlisting}

    \item \textbf{BiLSTM} (Bidirectional LSTM)
    \begin{lstlisting}[language=Python]
Sequential([
    Masking(mask_value=0.0),
    Bidirectional(LSTM(128, return_sequences=True)),
    Dropout(0.3),
    Bidirectional(LSTM(64)),
    Dropout(0.3),
    Dense(1, activation='sigmoid')
])
    \end{lstlisting}

    \item \textbf{GRU} (Gated Recurrent Unit)
    \begin{lstlisting}[language=Python]
Sequential([
    Masking(mask_value=0.0),
    GRU(128, return_sequences=True),
    Dropout(0.3),
    GRU(64),
    Dropout(0.3),
    Dense(1, activation='sigmoid')
])
    \end{lstlisting}
\end{enumerate}

\textbf{Training Configuration}:
\begin{itemize}
    \item Optimizer: Adam (lr=0.001)
    \item Loss: Binary crossentropy
    \item Batch size: 32
    \item Epochs: 50 (with early stopping, patience=10)
    \item Callbacks: EarlyStopping, ModelCheckpoint
\end{itemize}

\section{Experimental Results and Coaching Insights}

\subsection{Initial Classification Experiments}

As part of exploring whether automated stroke type identification could enable stroke-specific coaching, we trained multiple classification models:

\begin{table}[H]
\centering
\caption{Stroke Distinction Experiments (Test Set)}
\begin{tabular}{@{}lrrr@{}}
\toprule
\textbf{Model} & \textbf{Accuracy} & \textbf{F1 Score} & \textbf{ROC AUC} \\ \midrule
Random Forest & 0.424 & 0.344 & - \\
SVM & 0.479 & 0.521 & - \\
\midrule
LSTM & 0.457 & 0.452 & 0.467 \\
BiLSTM & 0.433 & 0.451 & 0.476 \\
GRU & 0.450 & 0.483 & 0.437 \\
\midrule
\textbf{Baseline (Random)} & \textbf{0.500} & \textbf{-} & \textbf{0.500} \\ \bottomrule
\end{tabular}
\end{table}

\textbf{Finding}: All models perform at random guessing level (50\% for balanced binary classification), indicating that these strokes cannot be reliably distinguished using pose features alone. This is an important insight for coaching system design.

\subsection{Cohen's d Effect Size Analysis}

Cohen's d measures the standardized difference between two groups:

\begin{equation}
d = \frac{\mu_{\text{Clear}} - \mu_{\text{Smash}}}{\sigma_{\text{pooled}}}
\end{equation}

where
\begin{equation}
\sigma_{\text{pooled}} = \sqrt{\frac{(n_1-1)\sigma_1^2 + (n_2-1)\sigma_2^2}{n_1+n_2-2}}
\end{equation}

\textbf{Interpretation}:
\begin{itemize}
    \item $|d| < 0.2$: Negligible effect
    \item $|d| = 0.2-0.5$: Small effect
    \item $|d| = 0.5-0.8$: Medium effect
    \item $|d| > 0.8$: Large effect
\end{itemize}

\subsubsection{Overall Feature Analysis}

\begin{table}[H]
\centering
\caption{Cohen's d Distribution Across 60 Features}
\begin{tabular}{@{}lr@{}}
\toprule
\textbf{Effect Size Category} & \textbf{Number of Features} \\ \midrule
Large ($|d| > 0.8$) & 0 \\
Medium ($|d| > 0.5$) & 0 \\
Small ($|d| > 0.2$) & 0 \\
Negligible ($|d| < 0.2$) & 60 \\ \bottomrule
\end{tabular}
\end{table}

\textbf{Conclusion}: No discriminative features found.

\subsubsection{Top 10 Features by Effect Size}

\begin{table}[H]
\centering
\caption{Most Discriminative Features (Still Negligible)}
\small
\begin{tabular}{@{}clrrrr@{}}
\toprule
\textbf{Rank} & \textbf{Feature} & \textbf{Cohen's d} & \textbf{Clear} & \textbf{Smash} & \textbf{p-value} \\ \midrule
1 & l\_forearm\_vertical\_angle\_vel & 0.141 & 0.014 & -0.094 & $<$0.001*** \\
2 & r\_forearm\_pitch\_vel & 0.130 & 0.029 & -0.070 & $<$0.001*** \\
3 & r\_forearm\_vertical\_angle\_vel & 0.130 & 0.029 & -0.070 & $<$0.001*** \\
4 & r\_wrist\_elbow\_vertical\_vel & -0.114 & 0.000 & 0.000 & $<$0.001*** \\
5 & l\_wrist\_elbow\_vertical\_vel & -0.092 & 0.000 & 0.000 & $<$0.001*** \\
6 & r\_arm\_plane\_pronation\_vel & -0.069 & 0.000 & 0.001 & $<$0.001*** \\
7 & r\_wrist\_horizontal\_reach & -0.042 & 0.074 & 0.076 & 0.013* \\
8 & r\_forearm\_vertical\_angle & \textbf{0.033} & \textbf{85.3°} & \textbf{84.2°} & 0.059 \\
9 & r\_forearm\_pitch & 0.033 & -4.7° & -5.8° & 0.059 \\
10 & r\_arm\_plane\_pronation & 0.030 & 0.027 & 0.008 & 0.081 \\ \bottomrule
\end{tabular}
\end{table}

\subsubsection{Wrist Orientation Features Analysis}

\textbf{Hypothesis}: Wrist angle differs significantly between Clear and Smash

\textbf{Result}: \textbf{Hypothesis REJECTED}

\begin{table}[H]
\centering
\caption{Wrist Feature Analysis}
\begin{tabular}{@{}lrrr@{}}
\toprule
\textbf{Feature} & \textbf{Cohen's d} & \textbf{Clear Mean} & \textbf{Smash Mean} \\ \midrule
r\_forearm\_vertical\_angle & \textbf{0.033} & \textbf{85.3°} & \textbf{84.2°} \\
r\_forearm\_pitch & 0.033 & -4.7° & -5.8° \\
r\_arm\_plane\_pronation & 0.030 & 0.027 & 0.008 \\
r\_wrist\_elbow\_vertical & -0.026 & 0.002 & 0.003 \\
r\_wrist\_horizontal\_reach & -0.042 & 0.074 & 0.076 \\
l\_forearm\_vertical\_angle & 0.009 & 82.9° & 82.7° \\ \bottomrule
\end{tabular}
\end{table}

\textbf{Key Finding}:
\begin{itemize}
    \item Expected difference: 60-90° (Clear upward, Smash downward)
    \item \textbf{Observed difference: 1.1°} (negligible)
    \item Cohen's d = 0.033 (far below 0.2 threshold)
\end{itemize}

\subsection{Visualization of Class Overlap}

Figure~\ref{fig:forearm_angle_dist} shows the distribution of forearm vertical angle for Clear and Smash strokes, demonstrating near-complete overlap.

\begin{figure}[H]
\centering
\begin{verbatim}
   Clear: ############################# (mean=85.3 deg, std=45.2 deg)
           ....###################....
                     ^
                   84-86 deg
                     v
  Smash: ############################ (mean=84.2 deg, std=44.8 deg)
           ....##################....

          0     30    60    90   120   150   180
                 Forearm Vertical Angle (degrees)
\end{verbatim}
\caption{Distribution of Forearm Vertical Angle (Key Feature)}
\label{fig:forearm_angle_dist}
\end{figure}

\section{Discussion: Implications for Coaching System Design}

\subsection{Biomechanical Insights from Analysis}

\subsubsection{Key Finding: Identical Body Mechanics}

Our analysis reveals an important coaching insight: at the moment of shuttlecock contact, Clear and Smash strokes exhibit virtually identical body poses:

\begin{table}[H]
\centering
\caption{Biomechanical Comparison at Contact Point}
\begin{tabular}{@{}lll@{}}
\toprule
\textbf{Characteristic} & \textbf{Clear} & \textbf{Smash} \\ \midrule
Racket position & Overhead & Overhead \\
Contact height & $\sim$2m above court & $\sim$2m above court \\
Arm extension & Fully extended & Fully extended \\
Body posture & Upright & Upright \\
Forearm angle & 85.3° $\pm$ 45.2° & 84.2° $\pm$ 44.8° \\
\midrule
\textbf{Difference} & \multicolumn{2}{c}{\textbf{1.1° (negligible)}} \\ \bottomrule
\end{tabular}
\end{table}

\subsubsection{Coaching Implication: Where Strokes Actually Differ}

From a coaching perspective, the key differences between Clear and Smash are \textbf{not in body positioning} but in racket control and outcome:

\begin{enumerate}
    \item \textbf{Racket Face Angle} (Critical coaching point)
    \begin{itemize}
        \item Clear: Open face, angled upward ($\sim$45° above horizontal)
        \item Smash: Closed face, angled downward ($\sim$30° below horizontal)
        \item \textbf{Coaching insight}: Players must focus on wrist control, not body positioning
        \item \textbf{Technical limitation}: Pose-only systems cannot capture racket orientation
    \end{itemize}

    \item \textbf{Shuttlecock Trajectory} (Outcome indicator)
    \begin{itemize}
        \item Clear: High parabolic arc (peak $\sim$6-8m)
        \item Smash: Steep downward line (angle $\sim$60-80° from horizontal)
        \item \textbf{Coaching insight}: Stroke effectiveness depends on outcome, not just form
    \end{itemize}

    \item \textbf{Follow-Through Motion} (Post-contact technique)
    \begin{itemize}
        \item Clear: Upward continuation, wrist extension
        \item Smash: Downward snap, wrist pronation
        \item \textbf{Future enhancement}: Extended temporal window needed
    \end{itemize}

    \item \textbf{Racket Head Speed} (Power generation)
    \begin{itemize}
        \item Clear: Moderate speed ($\sim$200-300 km/h)
        \item Smash: High speed ($\sim$300-400 km/h)
        \item \textbf{Coaching insight}: Body mechanics are similar; power comes from racket acceleration
    \end{itemize}
\end{enumerate}

\textbf{Key Coaching Takeaway}: Proper overhead stroke technique (body mechanics) is fundamentally the same for both Clear and Smash. The stroke type is determined by racket control at impact, not body positioning. This suggests coaching systems should focus on assessing the quality of fundamental overhead mechanics rather than trying to distinguish stroke types.

\subsection{Methodological Contributions}

Despite classification failure, this project demonstrates best practices:

\begin{enumerate}
    \item \textbf{Rigorous Statistical Validation}
    \begin{itemize}
        \item Cohen's d effect size analysis
        \item T-tests for significance
        \item Early detection of infeasible tasks
    \end{itemize}

    \item \textbf{Proper Data Handling}
    \begin{itemize}
        \item Group-stratified split (prevents leakage + maintains balance)
        \item Padding masking (proper sequence normalization)
        \item Global normalization (preserves class differences)
    \end{itemize}

    \item \textbf{Comprehensive Feature Engineering}
    \begin{itemize}
        \item Domain-driven features (wrist orientation)
        \item Multi-scale features (spatial + temporal)
        \item Systematic validation (Cohen's d per feature)
    \end{itemize}

    \item \textbf{Multiple Model Architectures}
    \begin{itemize}
        \item Traditional ML (RF, SVM) for baseline
        \item Deep learning (LSTM, BiLSTM, GRU) for temporal patterns
        \item Consistent results across architectures validates findings
    \end{itemize}
\end{enumerate}

\subsection{Limitations}

\subsubsection{Dataset Limitations}

\begin{itemize}
    \item \textbf{Temporal window}: Clips too short to capture follow-through
    \item \textbf{Missing modalities}: No racket tracking, no shuttlecock tracking
    \item \textbf{Camera variance}: Angles and lighting differ across matches
\end{itemize}

\subsubsection{MediaPipe Limitations}

\begin{itemize}
    \item \textbf{Body-only tracking}: 33 body landmarks, no hand details
    \item \textbf{2D projection}: Some 3D information lost
    \item \textbf{No equipment}: Cannot track racket or shuttlecock
\end{itemize}

\section{Revised Coaching System Approach}

\subsection{Pivot: From Stroke Classification to Technique Quality Assessment}

Based on our findings, we propose a revised coaching system architecture:

\subsubsection{Approach 1: Technique Quality Scoring (Recommended)}

\textbf{Concept}: Instead of distinguishing Clear from Smash, assess the quality of overhead stroke technique regardless of type.

\textbf{Method}:
\begin{itemize}
    \item Use existing pose features to define "good technique" indicators
    \item Train models to assess technique quality within each stroke type separately
    \item Provide quantitative feedback on key coaching points:
    \begin{itemize}
        \item Arm extension quality (0-100 score)
        \item Body posture alignment (0-100 score)
        \item Timing and balance (0-100 score)
        \item Swing smoothness (0-100 score)
    \end{itemize}
    \item Compare player's execution against expert benchmarks from dataset
\end{itemize}

\textbf{User Experience}:
\begin{enumerate}
    \item Player uploads video or manually selects stroke type
    \item System analyzes pose and computes technique scores
    \item Player receives specific feedback: "Your arm extension is 72/100. Try to reach higher at contact point."
\end{enumerate}

\subsubsection{Approach 2: Multi-Modal Enhancement (Future Work)}

For systems that require automatic stroke type detection, additional sensing is needed:

\textbf{Option A: Add Racket Tracking}
\begin{itemize}
    \item Train YOLO to detect badminton racket
    \item Compute racket face angle relative to court
    \item \textbf{Benefit}: Enables stroke type detection and racket control feedback
\end{itemize}

\textbf{Option B: Add Shuttlecock Tracking}
\begin{itemize}
    \item Track shuttlecock trajectory using computer vision
    \item Classify stroke type based on outcome trajectory
    \item \textbf{Benefit}: Direct measure of stroke effectiveness
\end{itemize}

\textbf{Option C: Extended Temporal Window}
\begin{itemize}
    \item Extend clips to include 1-2 seconds post-contact
    \item Analyze follow-through patterns
    \item \textbf{Benefit}: Capture full stroke mechanics including recovery
\end{itemize}

\subsection{Alternative: Pose-Based Classification for Different Stroke Pairs}

If stroke type classification remains desired, focus on strokes with distinct body mechanics:

\begin{table}[H]
\centering
\caption{Recommended Stroke Pairs for Pose-Based Coaching Systems}
\begin{tabular}{@{}lll@{}}
\toprule
\textbf{Stroke Pair} & \textbf{Expected d} & \textbf{Body Mechanics Difference} \\ \midrule
Overhead vs Underhand & $>$ 0.8 & Arm position (overhead vs waist) \\
Serve vs Return & $>$ 0.8 & Body motion (static vs dynamic) \\
Forehand vs Backhand & $>$ 0.8 & Arm side and rotation \\
Attacking vs Defensive & 0.5-0.8 & Stance and preparation \\ \bottomrule
\end{tabular}
\end{table}

\section{Conclusion and Next Steps}

This milestone report documents progress on developing an AI-based badminton coaching system for technique assessment. The work demonstrates rigorous methodology and provides valuable insights for coaching system design.

\subsection{Key Achievements}

\begin{enumerate}
    \item \textbf{Complete ML Pipeline}: Implemented end-to-end system from pose extraction to model evaluation
    \item \textbf{Domain-Driven Features}: Developed 60 biomechanical features including 9 novel wrist orientation features based on coaching requirements
    \item \textbf{Best Practices}: Group-stratified splitting, padding masking, proper sequence handling, statistical validation
    \item \textbf{Rigorous Analysis}: Cohen's d effect size analysis providing actionable coaching insights
    \item \textbf{System Architecture}: Established foundation for pose-based technique assessment
\end{enumerate}

\subsection{Critical Coaching Insight}

\textbf{Finding}: Clear and Smash strokes exhibit identical body mechanics at shuttlecock contact (Cohen's d $<$ 0.2). The stroke type is determined by racket control and outcome trajectory, not body positioning.

\textbf{Coaching Implication}: This validates that proper overhead technique is universal across stroke types. Players should focus on mastering fundamental overhead mechanics, then differentiate strokes through racket control.

\subsection{Value of This Work}

\begin{itemize}
    \item \textbf{For coaching}: Confirms that body mechanics coaching should be stroke-type agnostic
    \item \textbf{For system design}: Demonstrates that technique quality assessment is more appropriate than stroke classification for pose-based systems
    \item \textbf{For research}: Provides evidence-based guidance on feature discriminability in sports analytics
    \item \textbf{For responsible AI}: Shows importance of early validation to avoid wasted development effort
\end{itemize}

\subsection{Planned Next Steps}

Based on milestone findings, the project will pivot to:

\begin{enumerate}
    \item \textbf{Implement Technique Quality Scoring}
    \begin{itemize}
        \item Define quality metrics for overhead stroke technique
        \item Train regression models to score technique components
        \item Validate against expert benchmarks from dataset
    \end{itemize}

    \item \textbf{Develop Coaching Feedback System}
    \begin{itemize}
        \item Design user interface for video input and feedback display
        \item Implement quantitative scoring for arm extension, posture, timing
        \item Generate actionable recommendations for technique improvement
    \end{itemize}

    \item \textbf{Evaluation on User Need}
    \begin{itemize}
        \item Test system with sample stroke videos
        \item Validate that feedback aligns with coaching principles
        \item Assess system usability and practical value for players
    \end{itemize}

    \item \textbf{Documentation and Deployment}
    \begin{itemize}
        \item Implement demonstration interface (GUI or CLI)
        \item Document system capabilities and limitations
        \item Provide deployment instructions for end users
    \end{itemize}
\end{enumerate}

\textbf{Expected Final Report Contributions}:
\begin{itemize}
    \item Working coaching system with technique quality assessment
    \item Validated scoring methodology aligned with coaching principles
    \item Comprehensive evaluation demonstrating practical value
    \item Clear documentation of responsible AI considerations
    \item Recommendations for future multi-modal enhancements
\end{itemize}

\section*{Acknowledgments}

This project utilized the ShuttleSet dataset~\cite{ShuttleSet} and MediaPipe framework from Google. All code and analysis are available in the project repository.

\section*{References}

\begin{thebibliography}{1}
\bibitem{ShuttleSet}
Wei-Yao Wang, Yung-Chang Huang, Tsi-Ui Ik, and Wen-Chih Peng.
\textit{ShuttleSet: A Human-Annotated Stroke-Level Singles Dataset for Badminton Tactical Analysis}.
CoRR, abs/2306.04948, 2023.

\bibitem{MediaPipe}
Google Research.
\textit{MediaPipe Pose}.
\url{https://google.github.io/mediapipe/solutions/pose.html}, 2023.
\end{thebibliography}

\appendix

\section{Code Repository Structure}

\begin{verbatim}
iti123_v2/
|-- data/
|   |-- processed/
|   |   |-- poses/              # MediaPipe pose files
|   |   |-- features/           # Extracted features
|   |   +-- splits/             # Train/val/test splits
|-- src/
|   |-- data_processing/
|   |   |-- pose_extraction.py           # MediaPipe extraction
|   |   |-- feature_engineering_v2.py    # Feature extraction
|   |   +-- data_split.py                # Group-stratified split
|   |-- models/
|   |   |-- baseline_model.py            # RF, SVM
|   |   +-- lstm_model.py                # LSTM, BiLSTM, GRU
|   +-- analysis/
|       +-- analyze_wrist_features.py    # Cohen's d analysis
+-- outputs/
    +-- reports/
        |-- wrist_features_cohens_d.csv  # Full results
        +-- FINAL_PROJECT_REPORT.md      # Comprehensive report
\end{verbatim}

\section{Feature List}

\subsection{All 60 Sequence Features}

\begin{multicols}{2}
\begin{enumerate}[itemsep=0pt]
    \item r\_arm\_extension
    \item l\_arm\_extension
    \item arm\_extension\_ratio
    \item r\_wrist\_depth
    \item l\_wrist\_depth
    \item r\_wrist\_depth\_from\_hip
    \item r\_elbow\_depth
    \item body\_forward\_lean
    \item r\_wrist\_height\_rel
    \item l\_wrist\_height\_rel
    \item r\_wrist\_height\_from\_head
    \item r\_elbow\_height\_rel
    \item r\_wrist\_lateral
    \item shoulder\_width
    \item r\_elbow\_angle
    \item r\_shoulder\_angle
    \item torso\_lean
    \item shoulder\_rotation
    \item r\_upper\_arm\_length
    \item r\_forearm\_length
    \item arm\_straightness
    \item \textbf{r\_forearm\_vertical\_angle}
    \item \textbf{r\_forearm\_horizontal\_angle}
    \item \textbf{r\_wrist\_elbow\_vertical}
    \item \textbf{r\_wrist\_horizontal\_reach}
    \item \textbf{r\_arm\_plane\_pronation}
    \item \textbf{r\_arm\_plane\_vertical}
    \item \textbf{r\_forearm\_pitch}
    \item \textbf{l\_forearm\_vertical\_angle}
    \item \textbf{l\_wrist\_elbow\_vertical}
    \item [31-60] Velocity versions
\end{enumerate}
\end{multicols}

\textbf{Bold}: New wrist orientation features

\end{document}
